\scanpage{461}
\mspage{233}

Let us consider as reference element of $\mathcal H^{-}$ the element
\[
(40)\qquad \tau=\tau_{\varnothing}\qquad
\text{(anti-involution without fixed points of $X_g$)}
\]
we then find \inserted{by (32)} that for $I'\subset I$
\[
(41)\qquad \tau_{I'}=\rho_{I'}\tau=
\left(\prod_{j\in I'}\rho_j\right)\cdot\tau\quad
\deleted{\unsure{=\rho_h}}
\]
\margin{we see at once, since $\tau^2=1$, $\tau_{I'}^2=1$, that
$\tau_{I'}=\tau_{I'}^{-1}=\tau\rho_{I'}^{-1}$
\unsure{hence} $\tau\rho_{I'}\tau=\rho_{I'}^{-1}$
\unsure{which is $(42)$|which $\Rightarrow(42)$}}
moreover we will have, for \inserted{$\unsure{j|i}\in I$}
\[
\tau\rho_{i}\tau
=\tau_{\varnothing}\left(\tau_{\{\unsure{j|i}\}}\tau_{\{\varnothing\}}\right)
\tau_{\{0\}}
=\tau_{\varnothing}\tau_{\{j\}}=\rho_j^{-1}
\]
\[
(42)\qquad \tau\rho_{i}\tau^{-1}=\rho_{i}^{-1}.
\]

Thus $\mathcal G$ appears as the semidirect product of $\mathcal H\simeq\mathbb Z^I$, and of
\deleted{\illegible{2}}
$\{\pm1\}\simeq\{1,\tau\}$ operating on it by the symmetry
$\rho\longmapsto\rho^{-1}$. So for \uline{every} $\rho\in\mathcal H$,
\unsure{\deleted{$\tau\rho$} we have}
$(\tau\rho)^2=1$, i.e. for every
$\sigma\in\mathcal H^{-1}$,
$\sigma$ \deleted{\illegible{1}} \uline{is an anti-involution} !
Since $\sigma$ agrees with $\tau$ on
$V_0=X_g-\bigcup_{\unsure{j|i}\in I}\operatorname{Int}(T_j)$,
we see that the set of fixed points of $\sigma$ is contained in
$\bigcup_{i\in I}\operatorname{Int}(T_i)$, and to calculate the index of $\sigma$,
i.e. $\operatorname{card}\pi_0(X_g^\sigma)$, it suffices to take the sum of the
\margin{460}

\scanpage{462}

indices in the $T_j$. Now in $T_j$, we have \inserted{(29)}
$\rho_j(z,t)=(-z\exp\pi it,t)$, and by induction
\[
(43)\qquad \rho_j^{n_j}(z,t)=\left((-1)^{n_j}z\exp n_j\pi it,t\right)
\]
hence
\[
(44)\qquad \tau\rho_j^{n_j}(z,t)=
\left((-1)^{n_j+1}z\exp(n_j+1)\pi it,-t\right)
\]
and \deleted{\illegible{1}} $(z,t)$\inserted{$\in T_j$} is a fixed point of
$\unsure{z|\tau}\rho_j^{n_j}$ \unsure{iff|without}
$t=0$, \deleted{\unsure{if|and}} and $n_j$ is \uline{odd}, hence
\[
(45)\qquad T_j^{\tau\rho_j^{n_j}}=
\begin{cases}
\varnothing&\text{if $n_j$ even}\\
S_j\times\{0\}&\text{if $n_j$ odd}
\end{cases}
\]

hence
\[
\begin{gathered}
(46)\qquad \text{Index of }u=\tau\prod_j\rho_j^{n_j}\text{ is equal to the cardinality}\\
\text{of the set }I'_u\text{ of the }j\in I\text{ such that }n_j\\
\text{\unsure{is|are} odd}
\end{gathered}
\]

It follows \inserted{for general reasons} that $u$ is conjugate in
$\mathcal A_g=\operatorname{Aut}(X_g)$
\inserted{(by an element of $\unsure{A_g^+|\widetilde A_g^+}$)}
to $\tau_{I'}=\rho_{I'}\tau=\tau\rho_{I'}^{-1}$ \unsure{or also}
$\rho_{I'}^{-1}\tau=\tau\rho_{I'}$ ---
\deleted{
\unsure{but} \illegible{4}
\inserted{we have the equality \illegible{2} $\mathcal H$}
\unsure{in the group} $\mathcal G$ (\unsure{too small}).
\illegible{4} $\rho\in\mathcal H$, \illegible{2}
\unsure{the questions}, on the one hand
\[
\unsure{(\tau\rho)^{-1}}\cdot\rho_{I'}(\tau\rho)
=\unsure{\tau\rho\rho_{I'}}
\]
by introducing additional groups of symmetries, on the other hand by
\deleted{\unsure{enlarging it}, by} \unsure{completing}
each $T_j$ with a \unsure{configuration} more
}
\margin{\deleted{
\[
\unsure{(\tau\rho)(\tau\rho')(\tau\rho)^{-1}}
=\underset{\text{\deleted{\illegible{1}}}}{\tau\rho}\,
\underbrace{\tau\rho'\rho^{-1}\tau}_{\rho\rho'^{-1}}
=\tau\rho^{2}\rho'^{-1}=\tau\rho''
\]
as soon as $\rho^2=\rho'\rho''$, \unsure{i.e.|i.e.}
\illegible{3} $\tau\rho'$ is \unsure{conjugate} \illegible{1}
$\unsure{\tau\rho''}$ in $\mathcal G$, as soon as
$\rho'\rho''\in\mathcal H^2$
}}
\margin{461}

\scanpage{463}
\mspage{234}

But if $\tau\rho',\tau\rho''\in\mathcal G^-$, $(\rho',\rho''\in\mathcal H)$, and if
$\rho\in\mathcal H$, the relation
\[
\rho(\tau\rho')\rho^{-1}=\tau\rho''
\]
is equivalent to
\[
\underbrace{\tau\rho\tau}_{\rho^{-1}}\rho'\rho^{-1}=\rho''
\qquad\text{i.e.}\qquad \rho^2=\rho'\rho''^{-1}
\]
hence $\tau\rho'$ and $\tau\rho''$ are conjugate by an element of $\mathcal H$
\unsure{iff|without} $\rho'\rho''^{-1}\in\mathcal H^2$ -- which makes the preceding observation more precise\ldots

The way richest in simple symmetries for arranging the $2g$ antipodal holes
$D_j$ seems to me to be the following. We \deleted{\illegible{1}} consider
the Euclidean sphere, with \unsure{the} action of the dihedral group
$\unsure{\mathbb D_{2g}|\mathbf D_{2g}}$ -- for example, when it is the Riemann sphere that is
considered as Riemannian, by the choice of the antipodism as being
\[
(47)\qquad \tau z=-\frac{1}{\overline z}
\]
we take the standard action of the dihedral group (with the points $0,\infty$ as poles, and
the unit circle $\mathbb U=\{z\in\mathbb C\mathbin{:}|z|=1\}$ as equator),
\deleted{by} by writing
$\unsure{\mathbb D_n|\mathbf D_n}$ as $\subset\operatorname O(\deleted{\unsure{3}}2,\mathbb R)$,
as a $\frac12$ direct product of
\deleted{\unsure{$\mathbb Z/2\mathbb Z$}}\inserted{$\{\pm1\}$} by $\mu_n(\mathbb C)=\mu_n$,
the pair $(\zeta,\alpha)$ $(\zeta\in\mu_n,\alpha\in\deleted{\unsure{\mathbb Z/2\mathbb Z}}\inserted{\{\pm1\}})$ operating by
\[
(48)\qquad u_{\zeta,\alpha}(z)=\zeta z^\alpha
\]
\margin{462}

\scanpage{464}

We can moreover enlarge it to a group
$\mathbb D_n^\sim=\mathbb D_n\times\mathbb Z_{/2}$,
where the second factor $\mathbb Z_{/2}$ is generated by the antipodism (47)
(which commutes with $\mathbb D_n$ --- likewise moreover
does the anti-involution $z\longmapsto\frac{1}{\bar z}$, which has as its
set of fixed points $\mathbb U$ and is nothing other than the symmetry with respect to
the equator, their composite $z\longmapsto-z$ being the symmetry with respect to
the axis of the poles \ldots) hence to consider the transformations
\[
u_{\zeta,\alpha,\beta}\deleted{(z)}=\unsure{\zeta}u_{\zeta,\alpha}\tau^\beta
\qquad
(\zeta\in\mu_n,\ \alpha\in\{\pm1\},\ \beta\in\mathbb Z_{/2})
\]
hence
\[
u_{\zeta,\alpha,\beta}z=
\begin{cases}
u_{\zeta,\alpha}(z)=\zeta z^\alpha&\mathrm{si}\ \beta=0\\
u_{\zeta,\alpha}(-\frac{1}{\bar z})=-\zeta\bar z^{-\alpha}
&\mathrm{si}\ \beta=1
\end{cases},
\]
but we \unsure{will remember} that $\tau$ reverses the orientation,
hence \deleted{\illegible{1}}\inserted{is to be}
\deleted{\unsure{more}}\inserted{handled} with
caution as regards ``transport of structure'' in
the present situation.

Here $n=2g$, $\mathbb D_{\unsure{2g|n}}$ is of order $4g$, and
$\mathbb D_n^\sim$ of order $8g$. We take on the equator an
orbit $\unsure{\mathcal J|\mathbb J}$, of $\mu_n=\mu_{2g}$, for example
\unsure{the|here} \unsure{precisely}
\deleted{\unsure{the}} ($\mu_{2g}$
\inserted{\unsure{the set} of \unsure{$2g$th} roots of unity}
itself, \unsure{viewed}\deleted{\illegible{1}}
\unsure{as} a subset of the sphere) as the set of ``centers''
of the disks $D'_i$. We choose a disk $D'_0$ around the
\unsure{point|point} $P_0$ (small enough for what is to follow), and we
take the \unsure{transforms|transforms} of $D'_0$ by the
$\zeta\in\mu_n$.

\margin{it is \unsure{simply} necessary that $D'_0$ not
\unsure{meet} $\unsure{\zeta_0|\zeta}D'_0$,
\unsure{where} $\zeta=\exp\frac{2i\pi}{2g}$}

\margin{463}

\scanpage{465}
\mspage{235}

These choices having been made, the surface $X_g$ is determined without
\unsure{ambiguity|ambiguity}, and the group $\mathbb D_{2g}$ operates on it by transport of
structure, permuting the $g$ cylinders $T_i$ among themselves,
\unsure{corresponding|corresponding} to the elements of
\inserted{$\unsure{\mathcal J|\mathbb J}/\tau=$}$\unsure{\mathcal J|\mathbb J}/\{\pm1\}$ i.e.
to the pairs of antipodal elements of $\unsure{S|\mathcal J}$,
i.e. to the ``diagonals'' of the polygon with $2\mu$ sides
which they \deleted{\unsure{form}}\inserted{determine}
on the equator. The group $\mathbb D_{2g}$ normalizes
\deleted{\illegible{1}} the group $\mathcal G$, to be precise
we will have, for $u\in\mathbb D_{2g}$
\[
\tag{49}
u\tau_{I'}u^{-1}=\tau_{u(I')}
\]
for $I'\subset I=\unsure{\mathcal J|\mathbb J}/\tau$, taking into account
the operation of $\mathbb D_{2g}$ on $I$. We will thus have in
particular, denoting by $\tau_g$\inserted{$=\tau_\varnothing$}
the extension \margin{\deleted{(50)}} of the antipodism of the sphere $X_0$
(or rather of $\unsure{\tau}\vert V_0$) to an antipodism
without fixed points of $X_g$, denoted $\tau$ previously:
\[
\tag{50}
u\tau_g u^{-1}=\tau_g
\qquad\text{i.e. $\tau_g$ commutes with the action of $\mathbb D_n$}
\]
\[
\deleted{\unsure{(51)}}\qquad
u\rho_j u^{-1}=\rho_{u\unsure{(}j)}
\]

We can therefore say that $\mathbb D_{2g}$ operates on $\mathcal G$,
whence a $\frac12$-direct product $\mathbb D_{2g}.\mathcal G$,
which therefore operates on $X_g$.
As for the question of extending likewise
the action \inserted{on $X_0$,} of $\mathbb D_{2g}^\sim$
\deleted{on $X_0$} \unsure{as a whole} to an action on $X_g$,

\margin{464}

\scanpage{466}

this has been done without warning, to $\tau_{X_0}$ corresponding
naturally $\tau_{X_g}=\tau_g$, which indeed has the good taste
to commute with the action of $\mathbb D_{2g}$.
[It might seem more natural, it is true, in a
spirit of ``transport of structure (\unsure{in spite of} and against
everything?)'', to associate to $\tau_{X_0}$ the operation
$\tau_I$\inserted{\unsure{given|given} by (27)}, which on
each $T_j$ would be equal to $\tau'_j$
(and on $V_0$, of course, coincides with $\tau_{X_0}$),
$\tau'_j(z,t)=(z,-t)$,
\deleted{hence} the set of fixed points of $\tau_I^{\unsure{2}}$
being formed of the $g$ \unsure{middle} circles $S_j\times\{0\}$
of the $g$ tubes $T_j$. We will have by (41)
$\tau_I=\rho_I\tau_g$, hence $\tau_I$ commutes
\unsure{also} with $\mathbb D_n$, since $\tau$ and
$\rho_I=\prod_{i\in I}\rho_i$ commute with it.
But it \unsure{does|to me does} not seem \unsure{important} for the moment
which convention we \unsure{adopt}]

We can therefore also say that the group $\mathbb D_{2g}^\sim$
operates on $\mathcal H$ --- this operation extending that
of $\tau_g$\deleted{\unsure{by}}, now identified with an
element of $\mathbb D_{2g}^\sim$, i.e. with the antipodism in
$\mathbb D_{2g}^\sim$ --- and the $\frac12$ direct product
\[
\tag{51}
\mathbb D_{2g}^\sim.\mathcal H\supset\mathcal G
=\deleted{\unsure{\mathcal H}}\{1,\tau_g\}.\mathcal H
\]

\margin{465}